% Cases of test-notefoot.tex (\testN set) where footnotes at the foot meet the
% kernel's special output (review B1, B3), on letter paper, with \testN lines
% before the place it is about; \fnt{text} is a footnote with labels, as \fn's:
%   reins     a footnote goes to the foot (the margin is full); later on the
%             page a margin note's line goes to the next page (widow): the
%             kernel puts the foot back after that line;
%   reinsmix  the same with a footnote of LaTeX's own at the foot first;
%   wait      a footnote split at the foot and a second one in the same line,
%             which runs past the margin's foot or waits;
%   mcol3x    a footnote split at the foot and one that waits, then a paragraph:
%             the foot is counted, the text doesn't run over it;
%   mcol3     the same with multicols after them, with two footnotes in them: no
%             line of the columns over the foot;
%   mcol3n    the same as mcol3x with the two footnotes later on, in the margin
%             (in multicols with \testmulticols);
%   waitsec   the wait case after a heading and a line (shrinkable glue in the
%             page so far, review P2): the second footnote of the line stays
%             beside its line;
%   waitsh    the same with a \vspace{0pt minus 2pt} instead of the heading;
%   rest2     a footnote split at the foot, its rest heading the next page's
%             foot, where a footnote of LaTeX's own and a later margin note
%             whose line goes on (widow) follow, \testK lines down (review P1);
%   local     a footnote in a passage with \marginnoteoverflow{carry}, ending
%             before the paragraph does, where the margin is full: carried, as
%             the passage asks, though the mode is foot where the paragraph ends
%             (review S3); with \testlocalfoot, the other way round: the passage
%             asks for foot in a document with carry, and the note goes to the
%             foot;
%   defercarry  with \marginnotepriority{defer}: a deferrable note carried from
%             page 1 (its margin is full) gives way again on page 2 to a tall
%             note there, and a short footnote after it has room in the margin
%             of page 2 (review S5: counted at full height, the deferrable note
%             sent the footnote to the foot); with warn, the deferrable note
%             stays (as with off) and the footnote goes to the foot of page 2;
%   warnclaim with \marginnotepriority{warn}: a deferrable note, then a tall
%             note that doesn't fit with it, then a short footnote, on one page:
%             the deferrable note stays, so the footnote goes to the foot (it
%             was counted at nothing there, and the footnote went on with the
%             tall note);
%   queue     a tall margin note at the top of page 1, then \testQ floats for
%             the bottom: the first waits for the notes (said in the log), the
%             others wait behind it, as \LaTeX\ keeps floats in order; with
%             many, at the end on pages of floats;
%   cant      a chapter whose page comes out empty before it (as without
%             didactic): no warning that the callouts of a page can't be read,
%             as an empty column has none;
%   overmove  a short note after a tall one carried from the page: marginfix
%             keeps the notes in order, so it goes with the tall one, and it isn't
%             warned as running past the foot of the margin of a page it isn't on
%             (both are carried, each with its reason).
\makeatletter
\newcommand\fnt[1]{%
  \stepcounter{mk}\edef\mkv{\arabic{mk}}%
  \label{m-\mkv}\zsavepos{m-\mkv}%
  \footnote{\label{n-\mkv}\zsavepos{n-\mkv}#1}}
\newcommand\filler{Some words to fill up the rest of this line of the text and
  also quite a few of the next one, as the words of a paragraph do.}
\newcount\testli
\newcommand\testlines{\testli\z@
  \loop\ifnum\testli<\testN \noindent line \the\testli\par\advance\testli\@ne\repeat}
\def\testreins{reins}\def\testreinsmix{reinsmix}\def\testwait{wait}
\def\testmcolx{mcol3x}\def\testmcol{mcol3}\def\testmcoln{mcol3n}
\def\testwaitsec{waitsec}\def\testrestb{rest2}\def\testwaitsh{waitsh}
\def\testlocal{local}\def\testdefercarry{defercarry}\def\testwarnclaim{warnclaim}
\def\testqueue{queue}\def\testcant{cant}\def\testovermove{overmove}
\ifx\testcase\testmcolx\def\testfm#1{}\else\def\testfm#1{\fnt{#1}}\fi
\ifx\testcase\testwaitsh\def\testwaitsec{waitsh}\fi
\ifx\testcase\testwaitsec
  \noindent x\marginpar{\rule{1pt}{\dimexpr\textheight-8pt}} text.\par
  \testlines
  \ifx\testcase\testwaitsh\vspace{0pt minus 2pt}\else\section{SECX}\fi
  SECLINE one.\par
  CALLA\fnt{FOOTA \filler\ \filler\ \filler\ \filler\ \filler\ \filler\ ENDA}
  CALLW\fnt{FOOTW a second note, which waits for the next page.}
  \filler\ \filler\ \filler\ \filler\ \filler\ \filler\ \filler\ \filler\ ENDTEXT\par
\else\ifx\testcase\testrestb
  \widowpenalty=10000
  \noindent x\marginpar{\rule{1pt}{\dimexpr\textheight-8pt}} text.\par
  \testlines
  CALLA\fnt{FOOTA \filler\ \filler\ \filler\ \filler\ \filler\ \filler\ ENDA}
  \filler\ ENDTEXT\par
  {\footnotesatfoot \noindent PLAINCALL\footnote{FOOTPLAIN an ordinary footnote.}\par}
  \testli=0
  \loop\ifnum\testli<\testK \noindent pline \the\testli\par\advance\testli\@ne\repeat
  LINEONE\\ LINETWO\\ LINETHREE\testmp{NOTEB}\\ LINEFOUR\par
  After.
\else\ifx\testcase\testlocal
  \noindent x\marginpar{\rule{1pt}{\dimexpr\textheight-8pt}} text.\par
  \testlines
  \edef\testlocalmode{\ifdefined\testlocalfoot foot\else carry\fi}%
  {\expandafter\marginnoteoverflow\expandafter{\testlocalmode}%
  CALLL\fnt{FOOTL a footnote in a passage that asks for a mode of its own, with
  three lines or more of text, so that it can't run past the foot of the margin
  at all.}} more words.\par
  \filler\par
\else\ifx\testcase\testdefercarry
  \noindent FIRST\marginpar{\rule{1pt}{.5\textheight}} text.\par
  \testlines
  \noindent LTCALL\marginpar{\marginnotedeferrable LTNOTE\par
    \rule{1pt}{.6\textheight}} text.\par
  \testli=0
  \loop\ifnum\testli<18 \noindent more \the\testli\par\advance\testli\@ne\repeat
  \noindent PLAINCALL\marginpar{PNOTE\par\rule{1pt}{.6\textheight}} text.\par
  \noindent FNCALL\fnt{FOOTX a short footnote.} text.\par
  \filler\par
\else\ifx\testcase\testwarnclaim
  \noindent LTCALL\marginpar{\marginnotedeferrable LTNOTE\par
    \rule{1pt}{.6\textheight}} text.\par
  \testlines
  \noindent PLAINCALL\marginpar{PNOTE\par\rule{1pt}{.6\textheight}} text.\par
  \testlines
  \noindent FNCALL\fnt{FOOTX a short footnote.} text.\par
  \filler\par
\else\ifx\testcase\testqueue
  \noindent x\testmp{\rule{1pt}{.9\textheight}} text.\par
  \testlines
  \testli=0
  \loop\ifnum\testli<\testQ \advance\testli\@ne
    Words \the\testli.
    \begin{figure}[b]\centering\rule{4cm}{3\baselineskip}\caption{QF}\end{figure}%
  \repeat
  \par\filler\ \filler\ \filler\par
\else\ifx\testcase\testcant
  \chapter{Two}
  \noindent x\marginpar{\rule{1pt}{\dimexpr\textheight-8pt}} text.\par
  \vspace*{\dimexpr\textheight-6\baselineskip}
  \filler\ LOWCALL\fnt{FOOTLOW \filler\ \filler\ \filler\ \filler\ \filler\
  \filler\ \filler\ \filler\ \filler\ \filler\ \filler\ \filler}
  \filler\par
\else\ifx\testcase\testovermove
  \noindent x\testmp{\rule{1pt}{.5\textheight}} text.\par
  \noindent y\testmp{\rule{1pt}{.9\textheight}} text.\par
  \testlines
  \widowpenalty=10000
  LINEONE\\ LINETWO\testmp{OVNOTE}\\ LINETHREE\par
  \filler\par
\else\ifx\testcase\testwait
  \noindent x\marginpar{\rule{1pt}{\dimexpr\textheight-8pt}} text.\par
  \testlines
  CALLA\fnt{FOOTA \filler\ \filler\ \filler\ \filler\ \filler\ \filler\ ENDA}
  CALLW\fnt{FOOTW a second note, which waits for the next page.}
  \filler\ \filler\ \filler\ \filler\ \filler\ \filler\ \filler\ \filler\ ENDTEXT\par
\else\ifx\testcase\testreins\def\testkern{}\else\ifx\testcase\testreinsmix
  \def\testkern{{\footnotesatfoot KERN\footnote{KERNFN a footnote of LaTeX's own
    at the foot.}}}\fi\fi
\ifdefined\testkern
  \widowpenalty=10000
  \noindent x\marginpar{TALL\par\rule{1pt}{.97\textheight}} \testkern CALLA\fnt{FOOTA
  text of the footnote which is long enough to take several lines in the margin,
  more than two lines of text, so it cannot run past.} more text.\par
  \testlines
  LINEONE\\ LINETWO\\ LINETHREE\testmp{NOTEB}\\ LINEFOUR\par
  After.
\else
  \noindent x\marginpar{\rule{1pt}{\dimexpr\textheight-8pt}} text.\par
  \testlines
  \noindent CALLA\fnt{FOOTA \filler\ \filler\ \filler\ \filler\ \filler\ \filler\
  \filler\ \filler\ \filler\ \filler\ \filler\ \filler\ ENDA}CALLW\fnt{FOOTW the
  waiting note, with three lines or more of text so that it cannot run past the
  foot of the margin at all.}\par
  \ifx\testcase\testmcolx\else\ifdefined\testmulticols\begin{multicols}{2}\fi\fi
  \ifx\testcase\testmcol\begin{multicols}{2}\fi
  COLSTART\testfm{FOOTM in multicols.}\ \filler\ \filler\
  \filler\ \filler\ \filler\ \filler\ \filler\ \filler\ \filler\ \filler\ \filler\
  \filler\ \filler\ \filler\ \filler\ \filler\ \filler\ \filler\ \filler\ \filler\
  \filler\ \filler\ \filler\ \filler\ \filler\ \filler\ \filler\ \filler\ \filler\
  \filler\ COLLATE\testfm{FOOTM2 later in multicols.}\
  \filler\ \filler\ \filler\ \filler\ COLEND
  \ifx\testcase\testmcol\end{multicols}\fi
  \ifx\testcase\testmcolx\else\ifdefined\testmulticols\end{multicols}\fi\fi
  AFTER.
\fi\fi\fi\fi\fi\fi\fi\fi\fi\fi
\makeatother
